\documentclass[10pt,a4paper]{amsart}
\usepackage{amsmath,amssymb,amsthm,graphicx}
% For commutative diagrams we use Xy-pic
\usepackage[all]{xy}

% We use 2cell for 2-commutative diagrams.
\xyoption{2cell}
% Theorem environments.
%
\theoremstyle{plain}
\newtheorem{theorem}[subsection]{Theorem}
\newtheorem{proposition}[subsection]{Proposition}
\newtheorem{lemma}[subsection]{Lemma}

\theoremstyle{definition}
\newtheorem{definition}[subsection]{Definition}
\newtheorem{example}[subsection]{Example}
\newtheorem{exercise}[subsection]{Exercise}
\newtheorem{situation}[subsection]{Situation}

\theoremstyle{remark}
\newtheorem{remark}[subsection]{Remark}
\newtheorem{remarks}[subsection]{Remarks}

\numberwithin{equation}{subsection}

% Macros
%
\def\lim{\mathop{\mathrm{lim}}\nolimits}
\def\colim{\mathop{\mathrm{colim}}\nolimits}
\def\Spec{\mathop{\mathrm{Spec}}}
\def\Hom{\mathop{\mathrm{Hom}}\nolimits}
\def\Ext{\mathop{\mathrm{Ext}}\nolimits}
\def\SheafHom{\mathop{\mathcal{H}\!\mathit{om}}\nolimits}
\def\SheafExt{\mathop{\mathcal{E}\!\mathit{xt}}\nolimits}
\def\Sch{\mathit{Sch}}
\def\Mor{\mathop{\mathrm{Mor}}\nolimits}
\def\Ob{\mathop{\mathrm{Ob}}\nolimits}
\def\Sh{\mathop{\mathit{Sh}}\nolimits}
\def\NL{\mathop{N\!L}\nolimits}
\def\CH{\mathop{\mathrm{CH}}\nolimits}
\def\proetale{{pro\text{-}\acute{e}tale}}
\def\etale{{\acute{e}tale}}
\def\QCoh{\mathit{QCoh}}
\def\Ker{\mathop{\mathrm{Ker}}}
\def\Im{\mathop{\mathrm{Im}}}
\def\Coker{\mathop{\mathrm{Coker}}}
\def\Coim{\mathop{\mathrm{Coim}}}

% Boxtimes
%
\DeclareMathSymbol{\boxtimes}{\mathbin}{AMSa}{"02}

%
% Macros for moduli stacks/spaces
%
\def\QCohstack{\mathcal{QC}\!\mathit{oh}}
\def\Cohstack{\mathcal{C}\!\mathit{oh}}
\def\Spacesstack{\mathcal{S}\!\mathit{paces}}
\def\Quotfunctor{\mathrm{Quot}}
\def\Hilbfunctor{\mathrm{Hilb}}
\def\Curvesstack{\mathcal{C}\!\mathit{urves}}
\def\Polarizedstack{\mathcal{P}\!\mathit{olarized}}
\def\Complexesstack{\mathcal{C}\!\mathit{omplexes}}
% \Pic is the operator that assigns to X its picard group, usage \Pic(X)
% \Picardstack_{X/B} denotes the Picard stack of X over B
% \Picardfunctor_{X/B} denotes the Picard functor of X over B
\def\Pic{\mathop{\mathrm{Pic}}\nolimits}
\def\Picardstack{\mathcal{P}\!\mathit{ic}}
\def\Picardfunctor{\mathrm{Pic}}
\def\Deformationcategory{\mathcal{D}\!\mathit{ef}}
\usepackage{cleveref}
\crefname{originallemma}{original lemma}{original lemmas}
\Crefname{originallemma}{Original lemma}{Original lemmas}
\makeatletter
\newlabel{algebra-lemma-pushdown-module}{{02M0}{source}{}{sourcepointer-1}{}}
\newlabel{algebra-lemma-pushdown-module@cref}{{[originallemma][0][]02M0}{[1][0][]source}{}{}{}}
\makeatother
\crefname{originallemma}{original lemma}{original lemmas}
\Crefname{originallemma}{Original lemma}{Original lemmas}
\makeatletter
\newlabel{divisors-lemma-filter-after-modification}{{0AYP}{source}{}{sourcepointer-2}{}}
\newlabel{divisors-lemma-filter-after-modification@cref}{{[originallemma][0][]0AYP}{[1][0][]source}{}{}{}}
\makeatother
\crefname{originallemma}{original lemma}{original lemmas}
\Crefname{originallemma}{Original lemma}{Original lemmas}
\makeatletter
\newlabel{divisors-lemma-quasi-projective-Noetherian-pic-effective-Cartier}{{0AYM}{source}{}{sourcepointer-3}{}}
\newlabel{divisors-lemma-quasi-projective-Noetherian-pic-effective-Cartier@cref}{{[originallemma][0][]0AYM}{[1][0][]source}{}{}{}}
\makeatother
\crefname{originallemma}{original lemma}{original lemmas}
\Crefname{originallemma}{Original lemma}{Original lemmas}
\makeatletter
\newlabel{lemma-additivity-chern-classes}{{02UI}{source}{}{sourcepointer-4}{}}
\newlabel{lemma-additivity-chern-classes@cref}{{[originallemma][0][]02UI}{[1][0][]source}{}{}{}}
\makeatother
\crefname{originallemma}{original lemma}{original lemmas}
\Crefname{originallemma}{Original lemma}{Original lemmas}
\makeatletter
\newlabel{lemma-c1-cap-additive}{{02SP}{source}{}{sourcepointer-5}{}}
\newlabel{lemma-c1-cap-additive@cref}{{[originallemma][0][]02SP}{[1][0][]source}{}{}{}}
\makeatother
\crefname{originallemma}{original lemma}{original lemmas}
\Crefname{originallemma}{Original lemma}{Original lemmas}
\makeatletter
\newlabel{lemma-chern-classes-E-tensor-L}{{02UD}{source}{}{sourcepointer-6}{}}
\newlabel{lemma-chern-classes-E-tensor-L@cref}{{[originallemma][0][]02UD}{[1][0][]source}{}{}{}}
\makeatother
\crefname{originallemma}{original lemma}{original lemmas}
\Crefname{originallemma}{Original lemma}{Original lemmas}
\makeatletter
\newlabel{lemma-geometric-cap}{{02SQ}{source}{}{sourcepointer-7}{}}
\newlabel{lemma-geometric-cap@cref}{{[originallemma][0][]02SQ}{[1][0][]source}{}{}{}}
\makeatother
\crefname{originallemma}{original lemma}{original lemmas}
\Crefname{originallemma}{Original lemma}{Original lemmas}
\makeatletter
\newlabel{lemma-pushforward-cap-cj}{{02U9}{source}{}{sourcepointer-8}{}}
\newlabel{lemma-pushforward-cap-cj@cref}{{[originallemma][0][]02U9}{[1][0][]source}{}{}{}}
\makeatother
\crefname{originallemma}{original lemma}{original lemmas}
\Crefname{originallemma}{Original lemma}{Original lemmas}
\makeatletter
\newlabel{varieties-definition-degree-invertible-sheaf}{{0AYR}{source}{}{sourcepointer-9}{}}
\newlabel{varieties-definition-degree-invertible-sheaf@cref}{{[originallemma][0][]0AYR}{[1][0][]source}{}{}{}}
\makeatother
\crefname{originallemma}{original lemma}{original lemmas}
\Crefname{originallemma}{Original lemma}{Original lemmas}
\makeatletter
\newlabel{varieties-lemma-degree-additive}{{0AYS}{source}{}{sourcepointer-10}{}}
\newlabel{varieties-lemma-degree-additive@cref}{{[originallemma][0][]0AYS}{[1][0][]source}{}{}{}}
\makeatother
\crefname{originallemma}{original lemma}{original lemmas}
\Crefname{originallemma}{Original lemma}{Original lemmas}
\makeatletter
\newlabel{varieties-lemma-degree-birational-pullback}{{0AYU}{source}{}{sourcepointer-11}{}}
\newlabel{varieties-lemma-degree-birational-pullback@cref}{{[originallemma][0][]0AYU}{[1][0][]source}{}{}{}}
\makeatother
\crefname{originallemma}{original lemma}{original lemmas}
\Crefname{originallemma}{Original lemma}{Original lemmas}
\makeatletter
\newlabel{varieties-lemma-degree-effective-Cartier-divisor}{{0AYY}{source}{}{sourcepointer-12}{}}
\newlabel{varieties-lemma-degree-effective-Cartier-divisor@cref}{{[originallemma][0][]0AYY}{[1][0][]source}{}{}{}}
\makeatother
\crefname{originallemma}{original lemma}{original lemmas}
\Crefname{originallemma}{Original lemma}{Original lemmas}
\makeatletter
\newlabel{varieties-lemma-degree-in-terms-of-components}{{0AYW}{source}{}{sourcepointer-13}{}}
\newlabel{varieties-lemma-degree-in-terms-of-components@cref}{{[originallemma][0][]0AYW}{[1][0][]source}{}{}{}}
\makeatother
\crefname{originallemma}{original lemma}{original lemmas}
\Crefname{originallemma}{Original lemma}{Original lemmas}
\makeatletter
\newlabel{varieties-lemma-degree-tensor-product}{{0AYX}{source}{}{sourcepointer-14}{}}
\newlabel{varieties-lemma-degree-tensor-product@cref}{{[originallemma][0][]0AYX}{[1][0][]source}{}{}{}}
\makeatother
\crefname{originallemma}{original lemma}{original lemmas}
\Crefname{originallemma}{Original lemma}{Original lemmas}
\makeatletter
\newlabel{varieties-lemma-dim-1-proper-projective}{{0A26}{source}{}{sourcepointer-15}{}}
\newlabel{varieties-lemma-dim-1-proper-projective@cref}{{[originallemma][0][]0A26}{[1][0][]source}{}{}{}}
\makeatother
\title{The degree of a finite locally free module of constant rank on a proper curve equals the degree of its first Chern class capped with the fundamental class}
\author{The Stacks Project Authors}
\date{stacks-project:chow.tex:local-snapshot-20261012; GNU FDL 1.2 or later}
\begin{document}
\maketitle
\noindent\textit{Editorial note.} All delivered mathematical text is the author's own, copied from the pinned Stacks Project chapter file chow.tex; no statement or proof step was rewritten, reordered or removed.\par\medskip
\noindent\textit{Editorial note.} Preamble provenance: the pinned chapter file opens with \textbackslash{}input\{preamble\} and carries no preamble text of its own, so the project's single Stacks-wide preamble staged at sources/stacks-chow/preamble.tex is excerpted verbatim as two spans, lines 14-18 (Xy-pic setup) and lines 154-217 (the native theorem declarations and mathematical macros). Nothing was added, removed or redefined; the wrapper class and the amsmath/amssymb/amsthm/graphicx base packages are supplied by the assembler.\par\medskip
\noindent\textit{Editorial note.} Same-chapter pointers are located at their \textbackslash{}label lines inside the pinned chapter file; cross-chapter pointers are located in the project's own tag list staged in the snapshot (sources/stacks-chow/tags.txt), at the verbatim lines carrying their canonical tags.\par\medskip
\noindent\textit{Allowed standard prerequisites (editorial).}\begin{itemize}
\item Degree of an invertible sheaf on a proper curve and its behaviour under components, modification and tensor product. Varieties, labels varieties-definition-degree-invertible-sheaf (statement, proof line 8107), varieties-lemma-degree-in-terms-of-components (proof line 8117), varieties-lemma-degree-birational-pullback (proof line 8131), varieties-lemma-degree-additive (proof line 8135), varieties-lemma-degree-tensor-product (proof line 8150), varieties-lemma-degree-effective-Cartier-divisor (proof line 8155) and varieties-lemma-dim-1-proper-projective (proof line 8148). Supplies the notion of degree being compared, its decomposition over the irreducible components of X, its invariance under a birational modification, its additivity in a filtration and in tensor products, and projectivity of a proper curve; none of them contains the comparison with c\_1 itself.
\item Filtering a finite locally free sheaf on a curve by invertible quotients after a modification. Divisors, labels divisors-lemma-filter-after-modification (proof line 8122) and divisors-lemma-quasi-projective-Noetherian-pic-effective-Cartier (proof line 8143). Gives the modification X' -> X over which f\textasciicircum{}*E has a filtration with invertible successive quotients, and expresses an invertible sheaf on a proper curve as O(D) tensor O(D')\textasciicircum{}\{-1\} for effective Cartier divisors.
\item Compatibility of first Chern classes with pushforward, additivity and the geometric case. Same chapter chow.tex, labels lemma-pushforward-cap-cj (proof line 8126), lemma-additivity-chern-classes (proof line 8136), lemma-c1-cap-additive and lemma-chern-classes-E-tensor-L (proof line 8152) and lemma-geometric-cap (proof line 8157); Algebra, label algebra-lemma-pushdown-module (proof line 8172). Moves the computation across the modification, splits it additively over the filtration, reduces the invertible case to O(D), identifies c\_1(O(D)) cap [X] with [D]\_0, and converts the resulting sum of residue degrees into the dimension of the space of global sections of O\_D.
\end{itemize}\medskip
\ifdefined\hypertarget\hypertarget{sourcepointer-1}{}\fi
\noindent\textit{Source reference.} Original lemma 02M0: sources/stacks-chow/tags.txt, source lines 3427--3427. Navigation pointer for \textbackslash{}ref\{algebra-lemma-pushdown-module\} used inside the counted span (line 8172); the argument that consumes it is retained in the span.\par\medskip
\ifdefined\hypertarget\hypertarget{sourcepointer-2}{}\fi
\noindent\textit{Source reference.} Original lemma 0AYP: sources/stacks-chow/tags.txt, source lines 13797--13797. Navigation pointer for \textbackslash{}ref\{divisors-lemma-filter-after-modification\} used inside the counted span (line 8122); the argument that consumes it is retained in the span.\par\medskip
\ifdefined\hypertarget\hypertarget{sourcepointer-3}{}\fi
\noindent\textit{Source reference.} Original lemma 0AYM: sources/stacks-chow/tags.txt, source lines 13795--13795. Navigation pointer for \textbackslash{}ref\{divisors-lemma-quasi-projective-Noetherian-pic-effective-Cartier\} used inside the counted span (line 8143); the argument that consumes it is retained in the span.\par\medskip
\ifdefined\hypertarget\hypertarget{sourcepointer-4}{}\fi
\noindent\textit{Source reference.} Original lemma 02UI: sources/stacks-chow/pinned-original-source.tex, source lines 7965--7965. Navigation pointer for \textbackslash{}ref\{lemma-additivity-chern-classes\} used inside the counted span (line 8136); the argument that consumes it is retained in the span.\par\medskip
\ifdefined\hypertarget\hypertarget{sourcepointer-5}{}\fi
\noindent\textit{Source reference.} Original lemma 02SP: sources/stacks-chow/pinned-original-source.tex, source lines 4330--4330. Navigation pointer for \textbackslash{}ref\{lemma-c1-cap-additive\} used inside the counted span (line 8152); the argument that consumes it is retained in the span.\par\medskip
\ifdefined\hypertarget\hypertarget{sourcepointer-6}{}\fi
\noindent\textit{Source reference.} Original lemma 02UD: sources/stacks-chow/pinned-original-source.tex, source lines 7755--7755. Navigation pointer for \textbackslash{}ref\{lemma-chern-classes-E-tensor-L\} used inside the counted span (line 8152); the argument that consumes it is retained in the span.\par\medskip
\ifdefined\hypertarget\hypertarget{sourcepointer-7}{}\fi
\noindent\textit{Source reference.} Original lemma 02SQ: sources/stacks-chow/pinned-original-source.tex, source lines 4432--4432. Navigation pointer for \textbackslash{}ref\{lemma-geometric-cap\} used inside the counted span (line 8157); the argument that consumes it is retained in the span.\par\medskip
\ifdefined\hypertarget\hypertarget{sourcepointer-8}{}\fi
\noindent\textit{Source reference.} Original lemma 02U9: sources/stacks-chow/pinned-original-source.tex, source lines 7394--7394. Navigation pointer for \textbackslash{}ref\{lemma-pushforward-cap-cj\} used inside the counted span (line 8126); the argument that consumes it is retained in the span.\par\medskip
\ifdefined\hypertarget\hypertarget{sourcepointer-9}{}\fi
\noindent\textit{Source reference.} Original lemma 0AYR: sources/stacks-chow/tags.txt, source lines 13799--13799. Navigation pointer for \textbackslash{}ref\{varieties-definition-degree-invertible-sheaf\} used inside the counted span (line 8107); the argument that consumes it is retained in the span.\par\medskip
\ifdefined\hypertarget\hypertarget{sourcepointer-10}{}\fi
\noindent\textit{Source reference.} Original lemma 0AYS: sources/stacks-chow/tags.txt, source lines 13800--13800. Navigation pointer for \textbackslash{}ref\{varieties-lemma-degree-additive\} used inside the counted span (line 8135); the argument that consumes it is retained in the span.\par\medskip
\ifdefined\hypertarget\hypertarget{sourcepointer-11}{}\fi
\noindent\textit{Source reference.} Original lemma 0AYU: sources/stacks-chow/tags.txt, source lines 13802--13802. Navigation pointer for \textbackslash{}ref\{varieties-lemma-degree-birational-pullback\} used inside the counted span (line 8131); the argument that consumes it is retained in the span.\par\medskip
\ifdefined\hypertarget\hypertarget{sourcepointer-12}{}\fi
\noindent\textit{Source reference.} Original lemma 0AYY: sources/stacks-chow/tags.txt, source lines 13806--13806. Navigation pointer for \textbackslash{}ref\{varieties-lemma-degree-effective-Cartier-divisor\} used inside the counted span (line 8155); the argument that consumes it is retained in the span.\par\medskip
\ifdefined\hypertarget\hypertarget{sourcepointer-13}{}\fi
\noindent\textit{Source reference.} Original lemma 0AYW: sources/stacks-chow/tags.txt, source lines 13804--13804. Navigation pointer for \textbackslash{}ref\{varieties-lemma-degree-in-terms-of-components\} used inside the counted span (line 8117); the argument that consumes it is retained in the span.\par\medskip
\ifdefined\hypertarget\hypertarget{sourcepointer-14}{}\fi
\noindent\textit{Source reference.} Original lemma 0AYX: sources/stacks-chow/tags.txt, source lines 13805--13805. Navigation pointer for \textbackslash{}ref\{varieties-lemma-degree-tensor-product\} used inside the counted span (line 8150); the argument that consumes it is retained in the span.\par\medskip
\ifdefined\hypertarget\hypertarget{sourcepointer-15}{}\fi
\noindent\textit{Source reference.} Original lemma 0A26: sources/stacks-chow/tags.txt, source lines 12694--12694. Navigation pointer for \textbackslash{}ref\{varieties-lemma-dim-1-proper-projective\} used inside the counted span (line 8148); the argument that consumes it is retained in the span.\par\medskip
\begin{lemma}
\label{lemma-degree-vector-bundle}
Let $k$ be a field. Let $X$ be a proper scheme over $k$ of dimension $\leq 1$.
Let $\mathcal{E}$ be a finite locally free $\mathcal{O}_X$-module of constant
rank. Then
$$
\deg(\mathcal{E}) = \deg(c_1(\mathcal{E}) \cap [X]_1)
$$
where the left hand side is defined in
Varieties, Definition \ref{varieties-definition-degree-invertible-sheaf}.
\end{lemma}

\begin{proof}
Let $C_i \subset X$, $i = 1, \ldots, t$ be the irreducible components
of dimension $1$ with reduced induced scheme structure and let $m_i$ be the
multiplicity of $C_i$ in $X$. Then $[X]_1 = \sum m_i[C_i]$ and
$c_1(\mathcal{E}) \cap [X]_1$ is the sum of the pushforwards of the cycles
$m_i c_1(\mathcal{E}|_{C_i}) \cap [C_i]$. Since we have a similar decomposition
of the degree of $\mathcal{E}$ by
Varieties, Lemma \ref{varieties-lemma-degree-in-terms-of-components}
it suffices to prove the lemma in case $X$ is a proper curve over $k$.

\medskip\noindent
Assume $X$ is a proper curve over $k$.
By Divisors, Lemma \ref{divisors-lemma-filter-after-modification}
there exists a modification $f : X' \to X$ such that $f^*\mathcal{E}$
has a filtration whose successive quotients are invertible
$\mathcal{O}_{X'}$-modules. Since $f_*[X']_1 = [X]_1$ we conclude
from Lemma \ref{lemma-pushforward-cap-cj} that
$$
\deg(c_1(\mathcal{E}) \cap [X]_1) = \deg(c_1(f^*\mathcal{E}) \cap [X']_1)
$$
Since we have a similar relationship for the degree by
Varieties, Lemma \ref{varieties-lemma-degree-birational-pullback}
we reduce to the case where $\mathcal{E}$ has a filtration whose
successive quotients are invertible $\mathcal{O}_X$-modules.
In this case, we may use additivity of the degree
(Varieties, Lemma \ref{varieties-lemma-degree-additive})
and of first Chern classes (Lemma \ref{lemma-additivity-chern-classes})
to reduce to the case discussed in the next paragraph.

\medskip\noindent
Assume $X$ is a proper curve over $k$ and $\mathcal{E}$ is an
invertible $\mathcal{O}_X$-module. By
Divisors, Lemma
\ref{divisors-lemma-quasi-projective-Noetherian-pic-effective-Cartier}
we see that $\mathcal{E}$ is isomorphic to
$\mathcal{O}_X(D) \otimes \mathcal{O}_X(D')^{\otimes -1}$
for some effective Cartier divisors $D, D'$ on $X$ (this also uses
that $X$ is projective, see
Varieties, Lemma \ref{varieties-lemma-dim-1-proper-projective} for example).
By additivity of degree under tensor product of invertible sheaves
(Varieties, Lemma \ref{varieties-lemma-degree-tensor-product})
and additivity of $c_1$ under tensor product of invertible sheaves
(Lemma \ref{lemma-c1-cap-additive} or \ref{lemma-chern-classes-E-tensor-L})
we reduce to the case $\mathcal{E} = \mathcal{O}_X(D)$.
In this case the left hand side gives $\deg(D)$
(Varieties, Lemma \ref{varieties-lemma-degree-effective-Cartier-divisor})
and the right hand side gives $\deg([D]_0)$ by
Lemma \ref{lemma-geometric-cap}.
Since
$$
[D]_0 = \sum\nolimits_{x \in D}
\text{length}_{\mathcal{O}_{X, x}}(\mathcal{O}_{D, x}) [x] =
\sum\nolimits_{x \in D}
\text{length}_{\mathcal{O}_{D, x}}(\mathcal{O}_{D, x}) [x]
$$
by definition, we see
$$
\deg([D]_0) = \sum\nolimits_{x \in D}
\text{length}_{\mathcal{O}_{D, x}}(\mathcal{O}_{D, x}) [\kappa(x) : k] =
\dim_k \Gamma(D, \mathcal{O}_D) = \deg(D)
$$
The penultimate equality by
Algebra, Lemma \ref{algebra-lemma-pushdown-module}
using that $D$ is affine.
\end{proof}

\end{document}
